% !PRINTED 243
\weberpage{243}
Sind $x_1$ und $x_2$ beide imaginär, aber nicht conjugirt, so ist das Product
\[
(x_1 - x_2)(x_1' - x_2')
\]
und sein Quadrat positiv.

Sind aber endlich $x_1$ und $x_2$ conjugirt imaginär, so ist ihre Differenz \gesperrt{rein imaginär} und deren Quadrat negativ.

Es kommen also in dem Product, durch das wir die Discriminante erklärt haben, genau so viel negative Factoren vor, als es Paare conjugirt imaginärer Wurzeln giebt, und wir schliessen daraus auf den wichtigen Fundamentalsatz:

\gesperrt{Ist $f(x) = 0$ eine reelle Gleichung ohne mehrfache Wurzeln, so ist die Discriminante positiv oder negativ, je nachdem die Anzahl der Paare conjugirt imaginärer Wurzeln gerade oder ungerade ist.}

Der Fall, wo nur reelle Wurzeln vorhanden sind, gehört zu denen, wo die Discriminante positiv ist.

Für die Fälle der quadratischen und cubischen Gleichungen, in denen nur ein Paar imaginärer Wurzeln auftreten kann, ist durch diesen Hauptsatz bereits die Unterscheidung der verschiedenen Fälle, die in Bezug auf die Wurzelrealität möglich sind, die Determination, vollendet.

\begin{center}
\S.\,77.

\gesperrt{Discussion der quadratischen und cubischen Gleichung.}
\end{center}

Für die quadratische Gleichung
\[
a_0 x^2 + a_1 x + a_2 = 0
\]
haben wir (\S.\,46)
\[
\begin{array}{l}
D = a_1^2 - 4a_0a_2 > 0 \quad \text{reelle Wurzeln.}\\
\phantom{D = a_1^2 - 4a_0a_2 } < 0 \quad \text{imaginäre Wurzeln.}
\end{array}
\]
Für die cubische Gleichung
\[
a_0 x^3 + a_1 x^2 + a_2 x + a_3 = 0
\]
ist
\[
D = a_1^2a_2^2 + 18a_0a_1a_2a_3 - 4a_0a_2^3 - 4a_1^3a_3 - 27a_0^2a_3^2,
\]
\[
\begin{array}{l}
D > 0 \quad \text{drei reelle Wurzeln,}\\
D < 0 \quad \text{eine reelle, zwei imaginäre Wurzeln.}
\end{array}
\]
